\documentclass{article}
\usepackage{graphicx} % Required for inserting images
\usepackage[utf8]{inputenc}
\usepackage[greek,english]{babel}
\usepackage{alphabeta}
\usepackage{amsmath}  % Το κυριότερο για εξισώσεις
\usepackage{amssymb}  % Για σύμβολα όπως το κενό σύνολο, ανήκει, κλπ
\usepackage{amsfonts} % Για γραμματοσειρές μαθηματικών

\title{Differential Equations Exercises}
\author{Angelos Karadimitriou}
\date{Ongoing project - New exercises will be added}

\begin{document}

\maketitle

\section{Exercises}
\subsection{Exercise 1:}
Solve the differential equation :
$$y' = -\frac{1}{x^2} - \frac{y}{x} + y^2 $$
if one solution of the equation is $y_1=\frac{1}{x}$. \\\\
\textbf{Solution:} \\
This is a \textbf{Ricatti Equation}. We consider the following equation:
$$y(x) = y_1(x) + z = \frac{1}{x} + z$$
Then we differentiate with respect to $x$ :
$$y'(x) = -\frac{1}{x^2} + z'$$
Now we substitute for the given equation :
\begin{align*}
   -\frac{1}{x^2}+z' &= -\frac{1}{x^2}-\frac{\frac{1}{x}+z}{x}+\left(\frac{1}{x}+z\right)^2  \iff \\
z' &= \frac{-\frac{1}{x}-\frac{zx}{x}}{x}+\left(\frac{1}{x^2}+\frac{2z}{x}+z^2\right)  \\
&= \frac{\frac{-1-zx}{x}}{x}+\left(\frac{1}{x^2}+\frac{2z}{x}+z^2\right)  \\
&= \frac{-1-zx}{x^2}+\frac{1}{x^2}+\frac{2zx}{x^2}+z^2  \\
&= \frac{-1-zx+1+2zx}{x^2}+z^2  \\
&= \frac{z}{x}+z^2  \\
\end{align*}
Let $u = \frac{1}{z} \iff z = \frac{1}{u}$ which is a \textbf{Bernoulli Equation} substitution : 
\begin{align*}
    \left(\frac{1}{u}\right)'-\frac{1}{ux} &= \left(\frac{1}{u}\right)^2 \\
    \left(\frac{1}{u}\right)' &= \frac{1}{u^2}+\frac{1}{ux}  \\
    -\frac{1}{u^2}u' &= \frac{1}{u^2}+\frac{1}{ux}  \\
    -u' &= 1+\frac{u}{x} \\
    u'+\frac{u}{x} &= -1 
\end{align*}
Now we use the integrating factor : 
$$μ(x) = e^{\int\frac{1}{x}dx} = e^{lnx} = x $$
We multiply by the integrating factor and determine $u(x)$ as follows :
\begin{align*}
    (μ(x)u(x))' &= μ(x)(-1) \\
    xu(x) &= -\int{x}\,dx \\
    u(x)x &= -\frac{x^2}{2}+c \\
    u(x) &= -\frac{x}{2}+\frac{c}{x}
\end{align*}
Therefore, the general solution is :
\begin{align*}
    y(x) &= \frac{1}{x}+z \\
    y(x) &= \frac{1}{x}+\frac{1}{u(x)} \\
    y(x) &= \frac{1}{x}+\frac{1}{-\frac{x}{2}+\frac{c}{x}} \\ 
    y(x) &= \frac{1}{x}+\frac{2x}{2c-x^2} ,\quad c \in \mathbb{R}
\end{align*}

\newpage

\subsection{Exercise 2:}
Solve the differential equation :
$$x^2y'+2xy-y^3 = 0$$ 
\textbf{Solution:} \\
We will use the \textbf{Bernoulli Transformation} with $r = 3$. \\
Let: $$u(x) = y^{1-r}(x) = y^{-2}(x) = \frac{1}{y^{2}(x)}$$ 
Then we differentiate with respect to x:
$$ u'(x) = -2y^{-3}(x)y'(x) = -\frac{2}{y^3(x)}y'(x) $$
Now we multiply the given differential equation by:
$$(1-3)y^{-3}(x) = -\frac{2}{y^3(x)}$$
and we get:
$$-\frac{2}{y^3(x)}y'(x)+\frac{2}{x}y(x)\left(-\frac{2}{y^3(x)}\right) = \frac{1}{x^2}y^3\left(-\frac{2}{y^3(x)}\right) = -\frac{2}{x^2}$$
We now have the following differential equation:
$$u'(x)-\frac{4}{x}u(x) = -\frac{2}{x^2}$$
By using the integrating factor:
$$μ(x) = e^{\int-\frac{4}{x}}\,dx = e^{-4\int\frac{1}{x}}\,dx = e^{-4lnx} = e^{lnx^{-4}} = x^{-4} = \frac{1}{x^4} $$
we solve:
\begin{align*}
    (μ(x)u(x))' &= μ(x)\left(-\frac{2}{x^2}\right) \\
    \left(\frac{1}{x^4}\frac{1}{y^2(x)}\right)' &= \frac{1}{x^4}\left(-\frac{2}{x^2}\right) \\
    \frac{1}{x^4y^2(x)} &= -2\int x^{-6} \\
    \frac{1}{x^4y^2(x)} &= \frac{-2x^{-5}}{-5}+c \\
    \frac{1}{y^2(x)} &= \frac{2}{5x}+cx^4 \\
    \frac{1}{y^2} &= \frac{2+5cx^5}{5x} \\
    y^2(x) &= \frac{5x}{2+5cx} \\
    y(x) &= \pm\sqrt{\frac{5x}{2+5cx^5}}
\end{align*}

\newpage

\subsection{Exercise 3:}
Find the solution to the Initial Value Problem:
\begin{align*}
    y'(x) &= e^{-y}(5-4x) \\
    y(1) &= 0
\end{align*}
 \textbf{Solution:} \\
 This is a \textbf{Separable} differential equation :
\begin{align*}
     y'(x) &= e^{-y}(5-4x) \\
     \frac{dy}{dx} &= e^{-y}(5-4x) \\
     \frac{1}{e^{-y}}\,dy &= dx(5-4x) \\
     e^y\,dy &= (5-4x)\,dx \\
     \int{e^y}\,dy &= \int(5-4x)\,dx \\
     e^y &= 5x-2x^2+c \\
     \ln{e^y} &= \ln{(5x-2x^2+c)} \\
     y(x) &= ln{(5x-2x^2+c)} \\
\end{align*}
which represents the general solution of the equation. \\ \\
We substitute the given initial value $y(1) = 0$ :
\begin{align*}
    y(1) &= 0 \\
    \ln(5\cdot1-2\cdot1+c) &= 0 \\
    \ln{(3+c)} &= 0 \\
    e^{\ln{(3+c)}} &= e^0 \\
    3+c &= 1 \\
    c &= -2
\end{align*}
and we get the particular solution :
$$y(x) = \ln{(5x-2x^2-2)}$$

\newpage

\subsection{Exercise 4:} 
Find the solutions $y = y(x)$ of the differential equation:
$$\cos x-x\sin x +y^2+2xyy'(x) = 0$$
\textbf{Solution:} \\
Considering that $y'(x) = \frac{dy}{dx} $ we get:
$$(\cos x-x\sin x+y^2)\,dx+2xy\,dy = 0$$
we check if the differential equation is exact:
\begin{align*}
    M(x,y) &= \cos x-x\sin x+y^2 \\
    M_{y}(x,y) &= 2y \\
    N(x,y) &= 2xy \\
    N_{x}(x,y) &= 2y     
\end{align*}
We understand that:
$M_{y}(x,y) = N_{x}(x,y)$, which means that the differential equation is exact. \\
Let:
\begin{align}
    \Psi_{x}(x,y) &= \cos x-x\sin x+y^2  \label{ψx} \\
    \Psi_{y}(x,y) &= 2xy \label{ψy}
\end{align}
We integrate equation \eqref{ψx} with respect to $x$: \\
$$\int(\cos x-x\sin x+y^2)\,dx = \int(\cos x-x\sin x)\,dx+\int y^2\,dx = x\cos x+xy^2+h(y)$$
We differentiate with respect to y:
$$\Psi_{y}(x,y) = 0+2xy+h_{y}(y) = 2xy+h_{y}(y)$$
Comparing this with \eqref{ψy}:
$$2xy+h_{y}(y) = 2xy \iff h_{y}(y) = 0 \iff h(y) = c_1$$
Therefore, the general solution is:
$$\Psi(x,y) = c \iff x\cos x+xy^2+c_1 = c \iff x\cos x+xy^2 = c$$ \\
where c is an arbitrary constant

\newpage

 \subsection{Exercise 5:} 
 Find the general solution of the differential equation:
 $$y''(t)-\frac{1}{t}y'(t) = t\cos t,\,\, t>0$$
 \textbf{Solution:} \\
 We solve the differential equation using the \textbf{reduction of order} method: \\
 Let: 
 $$w(t) = y'(t)$$ 
 then 
 $$w'(t) = y''(t)$$ 
 This substitution reduces the equation to a first-order differential equation:
 $$w'(t)-\frac{1}{t}w(t) = t\cos t$$
 We find the integrating factor :
 $$μ(t) =  e^{-\int\frac{1}{t}\,dt} = e^{-\ln t} = \frac{1}{t} $$ \\
 and multiply:
 $$(w(t)μ(t))' = μ(t)t\cos t \iff \frac{w(t)}{t} = \int\cos t\,dt \iff w(t) = t\sin t+c_1t$$
 Integrating $w(t)$ with respect to t to find $y(t)$:
 \begin{align*}
      y(t) &= \int w(t) \,dt \\
      &= \int t\sin t \,dt + \int c_1t \,dt\\
      &= \int t(-\cos t)' \,dt + \int c_1t \,dt \\
      &= t(-\cos t)-\int(t)'(-\cos t) \,dt +\frac{c_1t^2}{2}\\
      &=-t\cos t+\int\cos t \,dt + \frac{c_1t^2}{2}\\
      &=-t\cos t+\sin t + \frac{c_1t^2}{2}+c_2 
 \end{align*}
 where $c_1,c_2$ are arbitrary constants.

\newpage

\subsection{Exercise 6:}
Find the general solution of the differential equation:
$$y''(t)+y(t) = 0, \,\, t \in \mathbb{R}$$
around the point $t = 0$, using the power series method. \\ \\
\textbf{Solution:} \\
Using the power series method:
$$y(t) = \sum_{n = 0}^{\infty}a_n(t-t_0)^n = \sum_{n = 0}^{\infty}a_nt^n$$
First derivative:
$$y'(t) = \sum_{n = 1}^{\infty}na_nt^{n-1}$$
Second derivative:
$$y''(t) = \sum_{n = 2}^{\infty}n(n-1)a_nt^{n-2}$$
Now we substitute the second derivative into the given equation:
\begin{align*}
    y''(t)+y(t) &= 0 \\
    \sum_{n = 2}^{\infty}n(n-1)a_nt^{n-2}+\sum_{n = 0}^{\infty}a_nt^n &= 0\\
    \sum_{k = 0}^{\infty}(k+2)(k+1)a_{k+2}t^k+\sum_{k = 0}^{\infty}a_kt^k &= 0 \\
    \sum_{n = 0}^{\infty}(n-2)(n-1)a_{n+2}t^n+\sum_{n = 0}^{\infty}a_nt^n &= 0 \\
    \sum_{n = 0}^{\infty}t^n((n+2)(n+1)a_{n+2}+a_n) = 0
\end{align*}
We solve the equation:
$$a_n+(n+2)(n+1)a_{n+2} = 0 \iff a_{n+2} = -\frac{a_n}{(n+2)(n+1)}$$
For even terms (n = 2k-2): \\ \\
For $n = 0: \, a_2 = -\frac{a_0}{1\cdot2} $ \\ \\
For $n = 2: \, a_4 = -\frac{a_2}{3\cdot4} = \frac{a_0}{1\cdot2\cdot3\cdot4}$ \\ \\
For $n = 4: \, a_6 = -\frac{a_4}{5\cdot6} = \frac{a_2}{3\cdot4\cdot5\cdot6} = -\frac{a_0}{1\cdot2\cdot3\cdot4\cdot5\cdot6}$ \\ \\
By using mathematical induction we get:
\begin{align*}
    a_{2k} = (-1)^k\frac{a_0}{(2k)!} \tag{1} \label{even}
\end{align*}    
For odd terms (2k-1): \\ \\
For $n = 1: \, a_3 = -\frac{a_1}{2\cdot3}$ \\ \\
For $n = 3: \, a_5 = -\frac{a_3}{4\cdot5} = \frac{a_1}{2\cdot3\cdot4\cdot5\cdot6}$ \\ \\
For $n = 5: \, a_7 = -\frac{a_5}{6\cdot7} = \frac{a_3}{4\cdot5\cdot6\cdot7} = -\frac{a_1}{2\cdot3\cdot4\cdot5\cdot6\cdot7}$\\ \\ 
By using mathematical induction we get:
\begin{align*}
    a_{2k+1} = (-1)^k\frac{a_1}{(2k+1)!} \tag{2} \label{odd} 
\end{align*}
Using \eqref{even} and \eqref{odd} we find the general solution: \\
\begin{align*}
    y(t) &= \sum_{k = 0}^{\infty}a_{2k}t^{2k}+\sum_{k = 0}^{\infty}a_{2k+1}t^{2k+1} \\
    &= a_0\sum_{k = 0}^{\infty}\frac{(-1)^kt^{2k}}{(2k)!}+a_1\sum_{k = 0}^{\infty}\frac{(-1)^kt^{2k+1}}{(2k+1)!} \\
    &= a_0\cos t + a_1\sin t
\end{align*}

\newpage

\subsection{Exercise 7:}
Consider the following system of differential equations: \\
$$x'(t) = 
\begin{pmatrix}
    -1 & 0 & 0 & 0 \\
    0 & 0 & 2 & 0 \\
    0 & -2 & 0 & 0 \\
    0 & 0 & 0 & 2
\end{pmatrix} 
x(t)$$
Find a fundamental matrix of solutions and identify all bounded solutions. \\ \\
\textbf{Solution:} \\
Let:
\begin{align*}
    u_{1} =
    \begin{pmatrix}
        u_{11} & u_{12} & u_{13} & u_{14} 
    \end{pmatrix}^T \\
    u_{2} =
    \begin{pmatrix}
        u_{21} & u_{22} & u_{23} & u_{24}
    \end{pmatrix}^T \\
    u_{3} = 
    \begin{pmatrix}
        u_{31} & u_{32} & u_{33} & u_{34}
    \end{pmatrix}^T \\
    u_{4} =
    \begin{pmatrix}
        u_{41} & u_{42} & u_{43} & u_{44}
    \end{pmatrix}^T \\
\end{align*}
We calculate the determinant: \\
\begin{align*}
    det(A-\lambda I_4) &= 0 \\
    det
    \begin{pmatrix}
    -1-\lambda & 0 & 0 & 0 \\
    0 & 0-\lambda & 2 & 0 \\
    0 & -2 & 0-\lambda & 0 \\
    0 & 0 & 0 & 2-\lambda \\
\end{pmatrix} 
&= 0 \\
(-1-\lambda)
\begin{vmatrix}
    -\lambda & 2 & 0 \\
    -2 & -\lambda & 0 \\
    0 & 0 & 2-\lambda \\ 
\end{vmatrix}
&= 0 \\
(-1-\lambda)
\begin{pmatrix}
    -\lambda
    \begin{vmatrix}
        -\lambda & 0 \\
        0 & 2-\lambda \\
    \end{vmatrix}
    -2
    \begin{vmatrix}
        -2 & 0 \\
        0 & 2-\lambda \\
    \end{vmatrix}
\end{pmatrix}
&= 0 \\
(-1-\lambda)[\lambda^2(2-\lambda)+4(2-\lambda)] &= 0 \\
-(1+\lambda)[(2-\lambda)(\lambda^2+4)] &=0 \\
\end{align*}
and we find the eigenvalues: \\
$$\lambda_{1} = -1, \, \lambda_{2} = 2, \, \lambda_{3} = 2i, \, \lambda_{4} = -2i $$ \\
Now we calculate the eigenvectors: \\ \\
For $\lambda_{1} = -1$: \\
$$(A-\lambda_{1}I_{4})u_{1} = 0 \iff 
\begin{pmatrix}
     0 & 0 & 0 & 0 \\
     0 & 1 & 2 & 0 \\
     0 & -2 & 1 & 0 \\
     0 & 0 & 0 & 3 \\
\end{pmatrix}
\begin{pmatrix}
     u_{11} \\
     u_{12} \\
     u_{13} \\
     u_{14} \\
\end{pmatrix}
= 0 $$
By solving, we find: \\
$$u_{11} = 1, \, u_{12} = u_{13} = u_{14} = 0 $$ \\
Thus:
$$u_{1} = 
\begin{pmatrix}
    1 & 0 & 0 & 0
\end{pmatrix}$$
and:
$$x_1(t) = e^{\lambda_{1}t}u_{1} = e^{-t}
\begin{pmatrix}
    1 \\
    0 \\
    0 \\
    0 \\
\end{pmatrix}$$
For $\lambda_{2} = 2:$
$$(A-\lambda_{2}I_{4})u_{2} = 0 \iff 
\begin{pmatrix}
    -3 & 0 & 0 & 0 \\
    0 & -2 & 2 & 0 \\
    0 &-2 & -2 & 0 \\
    0 & 0 & 0 & 0 \\
\end{pmatrix}
\begin{pmatrix}
    u_{21} \\
    u_{22} \\
    u_{23} \\
    u_{24} \\
\end{pmatrix} 
= 0$$
By solving, we find:
$$u_{21} = u_{22} = u_{23} = 0, \, u_{24} = 1$$
Thus:
$$u_{2} = 
\begin{pmatrix}
    0 & 0 & 0 & 1
\end{pmatrix}$$
and:
$$x_{2}(t) = e^{\lambda_{2}t}u_{2} = e^{2t} 
\begin{pmatrix}
    0 \\
    0 \\
    0 \\
    1 \\
\end{pmatrix}$$
For $\lambda_{3} = 2i$: \\
$$(A-\lambda_{3}I_{4})u_{3} = 0 \iff 
\begin{pmatrix}
    -1-2i & 0 & 0 & 0 & \\
    0 & -2i & 2 & 0 \\
    0 & -2 & -2i & 0 \\
    0 & 0 & 0 & 2-2i 
\end{pmatrix}
\begin{pmatrix}
    u_{31} \\
    u_{32} \\
    u_{33} \\
    u_{34} \\
\end{pmatrix}
= 0 $$
By solving, we find:
$$u_{31} = u_{34} = 0, \, u_{32} = 1, \, u_{33} = i$$
Thus:
$$u_{3} = 
\begin{pmatrix}
    0 & 1 & i & 0
\end{pmatrix}$$
and:
$$w(t) = e^{2it} 
\begin{pmatrix}
    0 \\
    1 \\
    i \\
    0 \\
\end{pmatrix} =
(\cos(2t)+i\sin(2t))
\begin{pmatrix}
    0 \\
    1 \\
    i \\
    0 \\
\end{pmatrix} =
\begin{pmatrix}
    0 \\
    \cos(2t)+i\sin(2t) \\
    -\sin(2t)+i\cos(2t) \\
    0
\end{pmatrix}$$
Now we get:
$$x_{3}(t) = Re(w) = 
\begin{pmatrix}
    0 \\
    \cos(2t) \\
    -\sin(2t) \\
    0 \\
\end{pmatrix}$$
and 
$$x_{4}(t) = Im(w) =
\begin{pmatrix}
    0 \\
    \sin(2t) \\
    \cos(2t) \\
    0 \\
\end{pmatrix}$$
Now we get \textbf{the Fundamental Matrix}:
$$\Phi(t) = 
\begin{pmatrix}
    e^{-t} & 0 & 0 & 0 \\
    0 & 0 & \cos(2t) & \sin(2t) \\
    0 & 0 & -\sin(2t) & \cos(2t) \\
    0 & e^{2t} & 0 & 0 \\
\end{pmatrix}$$
and the general solution:
$$x(t) = c_1
\begin{pmatrix}
    e^{-t} \\
    0 \\
    0 \\
    0 \\
\end{pmatrix}
+c_2
\begin{pmatrix}
    0 \\
    0 \\
    0 \\
    e^{2t} \\
\end{pmatrix}
+c_3
\begin{pmatrix}
    0 \\
    \cos(2t) \\
    -\sin(2t) \\
    0 \\
\end{pmatrix}
+c_4 
\begin{pmatrix}
    0 \\
    \sin(2t) \\
    \cos(2t) \\
    0 \\
\end{pmatrix}$$
The functions $\cos(2t)$ and $\sin(2t)$ are bounded (trigonometric functions). \\
The function $e^{-t}$ is bounded $\forall \, t$ \\
The function $e^{2t}$ is not bounded: \\
$$e^{2t} \xrightarrow{t\xrightarrow{}\infty} \infty$$
and therefore, for the solution to be bounded, we must have: $c_2 = 0$

\end{document}